% Ch14 self-study -- "I do / You do" ladder for Zumdahl chapter 14.
% SIXTEEN ladders (64 questions) plus a 12-question mixed set = 76
% questions, seventeen figures. Deliberately the longest document in
% the series: the teacher asked for extra length, and Unit 8 is
% 11-15% of the exam -- second only to Unit 3.
%
% SCOPE: nothing skipped. TEXTBOOK-MAP maps 14.1-14.11 onto CED 8.1,
% 8.2, 8.3, 8.6 and 8.7. (CED 8.4/8.5 -- titrations and buffers -- come
% from chapter 15 and are deliberately NOT previewed here.)
%
% Every pH in this file was checked against BOTH the small-x
% approximation and the exact quadratic, and the numbers were chosen so
% the two agree at the reported precision. See the fact-check script.
\documentclass{apchem-worksheet}

\apcunitword{Chapter}
\apcunit{14}
\apcunittitle{Acids and Bases}
\apcdoctitle{Self-Study \textbullet\ Chapter 14, I~do / You~do}
\apcsubtitle{Zumdahl \S14.1--14.11 \textbullet\ PDF pp.\ 697--740}

\begin{document}
\apctitleblock

\begin{apcnote}[How to use these notes --- read this first]
Each skill is a \textbf{ladder}: a fully worked example in the
solid-framed box, then \textbf{four} YOUR TURN questions in the dashed
box --- same skill, new numbers, no help. Work all four before checking
the gray \emph{check:} line; a tracker tick needs all four right first
time. After Ladder 16 there is a \textbf{mixed set} of twelve further
questions that do not tell you which ladder they came from.

\textbf{This is the longest document in the series, on purpose.} Unit 8
is \textbf{11--15\%} of the exam, second only to Unit 3, and this
chapter is most of it. Sixteen ladders, 76 questions.

\textbf{What you must bring with you.} Chapter 13. A weak acid problem
\emph{is} an equilibrium problem --- same ICE table, same small-$x$
approximation, same 5\% check. If Ladders 9--10 of Chapter 13 are not
automatic, stop and go back; nothing here will stick otherwise.

\textbf{The one habit that decides this chapter.} Before any
calculation, ask: \textbf{is it strong or weak?} Strong means it comes
apart completely and the arithmetic is one line. Weak means an
equilibrium and an ICE table. Answering that question first prevents
most of the errors this chapter can produce.

\textbf{Not covered here.} Buffers and titration curves (CED 8.4--8.5)
are Chapter 15. Do not go looking for them.
\end{apcnote}

% =====================================================================
\section*{\sffamily Ladder 1 \textbullet\ Three definitions}
\zsec{14.1, 14.11}

Three models, each broader than the last. Each is useful; none is wrong.

\begin{center}
\begin{tikzpicture}[font=\sffamily\footnotesize]
  % Columns at 0.3 / 3.5 / 8.3.  At 2.9 / 7.0 the name "Bronsted-Lowry"
  % overprinted its own acid text, and "acid: accepts an electron pair"
  % overprinted the base text beside it.
  \node[font=\sffamily\small\bfseries] at (6.4,3.6)
    {each definition contains the one before it};
  \foreach \y/\nm/\ac/\ba in {%
      2.85/{Arrhenius}/{produces \ce{H+} in water}/{produces \ce{OH-}
        in water},
      2.05/{Br\o nsted--Lowry}/{\textbf{donates} a proton}/{\textbf{accepts}
        a proton},
      1.25/{Lewis}/{\textbf{accepts} an electron pair}/{\textbf{donates}
        an electron pair}} {
    \node[anchor=west,font=\sffamily\bfseries] at (0.3,\y) {\nm};
    \node[anchor=west] at (3.5,\y) {acid: \ac};
    \node[anchor=west] at (8.3,\y) {base: \ba};}
  \draw[apcgray] (0.2,3.3) -- (12.9,3.3);
  \draw[apcgray] (0.2,0.85) -- (12.9,0.85);
  \node[anchor=west] at (0.3,0.4)
    {narrowest at the top, broadest at the bottom --- \ce{BF3} is a
     Lewis acid but has no proton to donate};
\end{tikzpicture}
\end{center}

\begin{workedexample}[I do: classifying under each model]
\textbf{Which model do you need for each of these?}

\textbf{$\ce{HCl(aq) -> H+(aq) + Cl-(aq)}$.} Arrhenius is enough:
\ce{HCl} produces \ce{H+} in water. It is also Br\o nsted--Lowry (it
donates a proton) and Lewis. All three agree.

\textbf{$\ce{NH3(aq) + H2O(l) <=> NH4+(aq) + OH-(aq)}$.} \ce{NH3}
contains no \ce{OH-} at all, so Arrhenius cannot call it a base. But it
\emph{accepts} a proton from water, so Br\o nsted--Lowry can. This is
exactly why the broader definition was needed.

\textbf{$\ce{BF3 + NH3 -> F3B-NH3}$.} No proton moves anywhere, so
Br\o nsted--Lowry has nothing to say. \ce{NH3} donates its lone pair and
\ce{BF3}, which is electron-deficient at boron, accepts it. Only the
Lewis model covers this.

\textbf{The relationship to keep straight.} Every Br\o nsted acid is a
Lewis acid, but not every Lewis acid is a Br\o nsted acid. The models
nest; the later ones do not replace the earlier ones.

\textbf{For the exam.} Br\o nsted--Lowry is the \emph{only} model
assessed --- the CED's topic 4.8 exclusion says outright that
\emph{``Lewis acid--base concepts will not be assessed on the AP
Exam''}, and Arrhenius appears only as history. The Lewis model stays
in this ladder because Zumdahl teaches it and it is the only way to
understand species like \ce{BF3} and hydrated metal cations --- but if
you are triaging, Br\o nsted--Lowry is the definition to own.
\end{workedexample}

\begin{yourturn}[YOUR TURN 1 --- four questions]
\begin{enumerate}[leftmargin=*,itemsep=0.5em]
  \item Define a Br\o nsted--Lowry acid. \workspace[1.4cm]
  \item Define a Lewis base. \workspace[1.4cm]
  \item Why can Arrhenius not classify \ce{NH3} as a base?
        \workspace[1.8cm]
  \item Is every Lewis acid a Br\o nsted acid? Give an example to
        support your answer. \workspace[1.8cm]
\end{enumerate}
\selfcheck{(a) proton donor \quad (b) electron-pair donor \quad (c) it
contains no \ce{OH-} \quad (d) no --- \ce{BF3}}
\keyonly{\begin{apcanswer}(a) A \textbf{proton (\ce{H+}) donor}.
(b) An \textbf{electron-pair donor}. (c) Because the Arrhenius
definition requires a base to \textbf{produce \ce{OH-}}, and \ce{NH3}
contains no hydroxide. It generates \ce{OH-} only indirectly, by taking
a proton from water --- which the Br\o nsted--Lowry definition
describes and the Arrhenius one cannot. (d) \textbf{No.} \ce{BF3} is a
Lewis acid because boron has an empty orbital and accepts an electron
pair, but it has \textbf{no hydrogen to donate}, so it is not a
Br\o nsted acid. The reverse does hold: every Br\o nsted acid is a Lewis
acid.\end{apcanswer}}
\end{yourturn}

% =====================================================================
\section*{\sffamily Ladder 2 \textbullet\ Conjugate acid--base pairs}
\zsec{14.1}

Every Br\o nsted reaction moves one proton, and that creates
\textbf{two} pairs. A conjugate pair differs by exactly one \ce{H+}.

\begin{center}
\begin{tikzpicture}[font=\sffamily\footnotesize]
  \node[font=\sffamily\small] at (5.3,3.3)
    {$\ce{HF}\;+\;\ce{H2O}\;\rightleftharpoons\;\ce{F-}\;+\;\ce{H3O+}$};
  \node[font=\sffamily\footnotesize] at (2.05,2.85) {acid};
  \node[font=\sffamily\footnotesize] at (3.75,2.85) {base};
  \node[font=\sffamily\footnotesize] at (6.0,2.85) {conj.\ base};
  \node[font=\sffamily\footnotesize] at (8.1,2.85) {conj.\ acid};
  % pair 1: HF / F-
  \draw[thick] (2.05,2.55) -- (2.05,1.95) -- (6.0,1.95) -- (6.0,2.55);
  \node[font=\sffamily\footnotesize] at (4.0,1.65)
    {pair 1: differ by one \ce{H+}};
  % pair 2: H2O / H3O+
  \draw[thick] (3.75,2.55) -- (3.75,1.1) -- (8.1,1.1) -- (8.1,2.55);
  \node[font=\sffamily\footnotesize] at (5.9,0.8)
    {pair 2: differ by one \ce{H+}};
  \draw[apcgray] (0.3,0.35) -- (10.4,0.35);
  \node[anchor=west] at (0.4,-0.1)
    {\textbf{lose} a proton to get the conjugate base \textbullet\
     \textbf{gain} one to get the conjugate acid};
\end{tikzpicture}
\end{center}

\begin{apctrap}
\textbf{A conjugate pair differs by \emph{one} proton --- and by one
unit of charge.} \ce{H2SO4} and \ce{SO4^2-} are \textbf{not} a
conjugate pair: they differ by two protons. The conjugate base of
\ce{H2SO4} is \ce{HSO4-}. Check the charge as well as the hydrogen
count and you will not get this wrong.
\end{apctrap}

\begin{workedexample}[I do: finding conjugates]
\textbf{Rule: to get a conjugate \emph{base}, remove one \ce{H+}
(charge goes down by one). To get a conjugate \emph{acid}, add one
\ce{H+} (charge goes up by one).}

\textbf{Conjugate base of \ce{HNO3}:} remove \ce{H+} from \ce{HNO3}
(charge 0) to get \ce{NO3-} (charge $-1$). \checkmark

\textbf{Conjugate acid of \ce{NH3}:} add \ce{H+} to \ce{NH3} (charge 0)
to get \ce{NH4+} (charge $+1$). \checkmark

\textbf{Conjugate base of \ce{H2PO4^-}:} remove \ce{H+} from charge
$-1$ to get \ce{HPO4^2-}, charge $-2$. \checkmark

\textbf{Now the interesting one --- \ce{HCO3-}.} Its conjugate base is
\ce{CO3^2-} and its conjugate acid is \ce{H2CO3}. A species that can do
both is \textbf{amphoteric}, and water is the most important example:
\ce{H2O} has conjugate base \ce{OH-} and conjugate acid \ce{H3O+}.

\textbf{The pattern worth internalising.} Strong acid, weak conjugate
base. \ce{HCl} is a strong acid, so \ce{Cl-} is such a weak base that it
does nothing in water. Weak acid, appreciable conjugate base: acetic
acid is weak, so acetate is a genuine base --- which is exactly why
sodium acetate solutions are basic (Ladder 14).
\end{workedexample}

\begin{yourturn}[YOUR TURN 2 --- four questions]
\begin{enumerate}[leftmargin=*,itemsep=0.5em]
  \item Give the conjugate base of \ce{HCN}. \workspace[1.4cm]
  \item Give the conjugate acid of \ce{HPO4^2-}. \workspace[1.4cm]
  \item In $\ce{HCl + H2O -> Cl- + H3O+}$, name both conjugate pairs.
        \workspace[1.8cm]
  \item Why are \ce{H2SO4} and \ce{SO4^2-} not a conjugate pair?
        \workspace[1.6cm]
\end{enumerate}
\selfcheck{(a) \ce{CN-} \quad (b) \ce{H2PO4^-} \quad (c) \ce{HCl}/
\ce{Cl-} and \ce{H2O}/\ce{H3O+} \quad (d) they differ by two protons}
\keyonly{\begin{apcanswer}(a) \textbf{\ce{CN-}}. (b)
\textbf{\ce{H2PO4^-}} --- add one \ce{H+}, charge rises from $-2$ to
$-1$. (c) \textbf{\ce{HCl}/\ce{Cl-}} and
\textbf{\ce{H2O}/\ce{H3O+}}. (d) Because they differ by \textbf{two}
protons, not one. The conjugate base of \ce{H2SO4} is \ce{HSO4-}, and
the conjugate base of \emph{that} is \ce{SO4^2-}.\end{apcanswer}}
\end{yourturn}

% =====================================================================
\section*{\sffamily Ladder 3 \textbullet\ Strong versus weak}
\zsec{14.2}

``Strong'' has one meaning only: it ionises \textbf{completely}. It has
nothing to do with concentration.

\begin{center}
\begin{tikzpicture}[font=\sffamily\footnotesize]
  % strong: all dissociated
  \node[font=\sffamily\small\bfseries] at (2.2,3.4) {strong acid};
  \fill[apcgray!22] (0.6,0.5) rectangle (3.8,2.9);
  \draw[very thick] (0.6,2.9) -- (0.6,0.5) -- (3.8,0.5) -- (3.8,2.9);
  \foreach \x/\y in {1.0/0.9, 1.7/1.3, 2.5/0.95, 3.3/1.35,
                     1.2/1.9, 2.1/2.3, 2.9/1.85, 3.5/2.35} {
    \node[font=\tiny] at (\x,\y) {\ce{H+}};}
  \node[align=center] at (2.2,-0.15)
    {\textbf{no} intact \ce{HA} left\\one-way arrow};

  % weak: mostly intact
  \begin{scope}[xshift=6.0cm]
    \node[font=\sffamily\small\bfseries] at (2.2,3.4) {weak acid};
    \fill[apcgray!22] (0.6,0.5) rectangle (3.8,2.9);
    \draw[very thick] (0.6,2.9) -- (0.6,0.5) -- (3.8,0.5) -- (3.8,2.9);
    \foreach \x/\y in {1.0/0.9, 1.9/1.25, 2.8/0.95, 3.4/1.4,
                       1.3/1.9, 2.3/2.3, 3.2/1.9} {
      \node[font=\tiny] at (\x,\y) {\ce{HA}};}
    \node[font=\tiny] at (2.0/1,2.75) {\ce{H+}};
    \node[align=center] at (2.2,-0.15)
      {\textbf{mostly} intact \ce{HA}\\equilibrium arrow};
  \end{scope}
\end{tikzpicture}
\end{center}

\begin{center}
\begin{tikzpicture}[font=\sffamily\footnotesize]
  \node[font=\sffamily\small\bfseries] at (5.3,2.55)
    {learn these lists --- anything not on them is weak};
  % list column at 3.9, not 3.3: "strong acids (6):" in bold sans is
  % 2.9cm wide and ran straight into the first formula.
  \node[anchor=west,font=\sffamily\bfseries] at (0.4,1.95)
    {strong acids (6):};
  \node[anchor=west] at (3.9,1.95)
    {\ce{HCl}, \ce{HBr}, \ce{HI}, \ce{HNO3}, \ce{H2SO4}, \ce{HClO4}};
  \node[anchor=west,font=\sffamily\bfseries] at (0.4,1.35)
    {strong bases:};
  \node[anchor=west] at (3.9,1.35)
    {group 1 hydroxides, plus \ce{Ca(OH)2}, \ce{Sr(OH)2},
     \ce{Ba(OH)2}};
  \draw[apcgray] (0.3,0.95) -- (10.5,0.95);
  \node[anchor=west] at (0.4,0.5)
    {note \ce{HF} is \textbf{not} strong --- it is the one halide acid
     that is weak (Ladder 15 explains why)};
\end{tikzpicture}
\end{center}

\begin{apctrap}
\textbf{Strong is not the same as concentrated.} A
\SI{0.001}{\Molar} \ce{HCl} solution is \emph{dilute} but the acid is
\emph{strong}; a \SI{10}{\Molar} acetic acid solution is
\emph{concentrated} but the acid is \emph{weak}. Strong/weak describes
the substance and how completely it ionises; dilute/concentrated
describes how much of it you put in the water. The two words answer
different questions.
\end{apctrap}

\begin{workedexample}[I do: what ``complete'' actually buys you]
\textbf{Compare \SI{0.10}{\Molar} \ce{HCl} with \SI{0.10}{\Molar}
\ce{CH3COOH}.}

\textbf{\ce{HCl} is strong.} It ionises completely, so
\[ \ce{HCl -> H+ + Cl-} \]
is written with a \emph{one-way} arrow, and $[\ce{H+}] =
\SI{0.10}{\Molar}$ exactly. There is no equilibrium and no ICE table:
whatever you put in, you get out as ions.

\textbf{Acetic acid is weak.} It reaches an equilibrium,
\[ \ce{CH3COOH <=> H+ + CH3COO-} \]
written with a \emph{double} arrow, and only about 1\% of it comes
apart. $[\ce{H+}]$ is around \SI{1.3e-3}{\Molar} --- about
seventy-five times smaller than for the \ce{HCl}, at the same nominal
concentration.

\textbf{So the arrow you draw is a claim.} A single arrow asserts
complete ionisation; a double arrow asserts an equilibrium. Drawing the
wrong one is a chemistry error, not a notation slip, and examiners treat
it as such.

\textbf{And the practical consequence.} For a strong acid you can go
straight to pH in one step (Ladder 6). For a weak acid you must build an
ICE table and use $K_{a}$ (Ladder 9). Deciding strong-or-weak is
therefore always your first move, before any arithmetic at all.
\end{workedexample}

\begin{yourturn}[YOUR TURN 3 --- four questions]
\begin{enumerate}[leftmargin=*,itemsep=0.5em]
  \item Is \ce{HF} a strong or a weak acid? \workspace[1.4cm]
  \item Which is more acidic: \SI{0.10}{\Molar} \ce{HCl} or
        \SI{0.10}{\Molar} \ce{HF}? Why? \workspace[1.8cm]
  \item Explain the difference between ``strong'' and
        ``concentrated''. \workspace[1.8cm]
  \item Which arrow, single or double, do you draw for
        \ce{HNO3} in water? Why? \workspace[1.6cm]
\end{enumerate}
\selfcheck{(a) weak \quad (b) \ce{HCl} --- it ionises completely \quad
(c) extent of ionisation vs amount dissolved \quad (d) single ---
\ce{HNO3} is strong}
\keyonly{\begin{apcanswer}(a) \textbf{Weak} --- the only hydrohalic acid
that is. (b) \textbf{\ce{HCl}}. It is strong and ionises completely, so
$[\ce{H+}] = \SI{0.10}{\Molar}$; \ce{HF} is weak and only partly
ionises, giving a much smaller $[\ce{H+}]$ at the same nominal
concentration. (c) \textbf{Strong} describes \emph{how completely the
substance ionises} --- a property of the acid itself.
\textbf{Concentrated} describes \emph{how much was dissolved} --- a
property of the solution. A dilute solution of a strong acid and a
concentrated solution of a weak acid are both perfectly ordinary. (d)
A \textbf{single} arrow, because \ce{HNO3} is one of the six strong
acids and ionisation is complete --- there is no equilibrium to
represent.\end{apcanswer}}
\end{yourturn}

% =====================================================================
\section*{\sffamily Ladder 4 \textbullet\ Water, $K_{w}$, and
autoionisation}
\zsec{14.2--14.3}

Water ionises slightly on its own, and that single equilibrium underlies
the entire pH scale.

\[ \ce{2H2O(l) <=> H3O+(aq) + OH-(aq)}
   \qquad
   K_{w} = [\ce{H+}][\ce{OH-}] = \num{1.0e-14}
   \text{ at \SI{25}{\celsius}} \]

\begin{center}
\begin{tikzpicture}[font=\sffamily\footnotesize]
  \node[font=\sffamily\small\bfseries] at (5.3,3.3)
    {the four quantities, and the three bridges between them};
  \node[font=\sffamily\bfseries] at (1.6,2.5) {$[\ce{H+}]$};
  \node[font=\sffamily\bfseries] at (9.0,2.5) {$[\ce{OH-}]$};
  \node[font=\sffamily\bfseries] at (1.6,1.0) {pH};
  \node[font=\sffamily\bfseries] at (9.0,1.0) {pOH};
  \draw[{Stealth[length=2mm]}-{Stealth[length=2mm]},thick]
    (2.4,2.5) -- (8.2,2.5);
  \node[fill=white,font=\sffamily\footnotesize] at (5.3,2.5)
    {$[\ce{H+}][\ce{OH-}] = \num{1.0e-14}$};
  \draw[{Stealth[length=2mm]}-{Stealth[length=2mm]},thick]
    (2.4,1.0) -- (8.2,1.0);
  \node[fill=white,font=\sffamily\footnotesize] at (5.3,1.0)
    {pH $+$ pOH $= 14.00$};
  \draw[{Stealth[length=2mm]}-{Stealth[length=2mm]},thick]
    (1.6,2.15) -- (1.6,1.35);
  \node[anchor=east,font=\sffamily\footnotesize] at (1.35,1.75)
    {$-\log$};
  \draw[{Stealth[length=2mm]}-{Stealth[length=2mm]},thick]
    (9.0,2.15) -- (9.0,1.35);
  \node[anchor=west,font=\sffamily\footnotesize] at (9.25,1.75)
    {$-\log$};
  \node[font=\sffamily\footnotesize] at (5.3,0.35)
    {know any one and you can reach the other three};
\end{tikzpicture}
\end{center}

\begin{apctrap}
\textbf{``pH 7 is neutral'' is only true at \SI{25}{\celsius}.} $K_{w}$
grows with temperature, so at \SI{50}{\celsius} neutral water has
pH $\approx 6.6$ --- and it is still \emph{neutral}, because
$[\ce{H+}]$ still equals $[\ce{OH-}]$. Neutral means the two are equal,
not that pH is 7. Any question that gives you a temperature other than
\SI{25}{\celsius} is testing exactly this.
\end{apctrap}

\begin{workedexample}[I do: moving around the wheel]
\textbf{A solution has $[\ce{H+}] = \SI{2.0e-4}{\Molar}$ at
\SI{25}{\celsius}. Find $[\ce{OH-}]$, pH and pOH.}

\textbf{$[\ce{OH-}]$ from $K_{w}$:}
\[ [\ce{OH-}] = \frac{\num{1.0e-14}}{\num{2.0e-4}}
   = \SI{5.0e-11}{\Molar} \]

\textbf{pH by definition:}
\[ \text{pH} = -\log(\num{2.0e-4}) = \mathbf{3.70} \]

\textbf{pOH two ways, and they must agree:}
\[ \text{pOH} = 14.00 - 3.70 = \mathbf{10.30}
   \qquad\text{or}\qquad
   -\log(\num{5.0e-11}) = 10.30 \quad\checkmark \]

\textbf{Is the answer sensible?} pH 3.70 is below 7, so the solution is
acidic --- consistent with $[\ce{H+}]$ being far larger than
$[\ce{OH-}]$. That one-line check catches an inverted $K_{w}$ division,
which is the usual slip here.

\textbf{Note about significant figures in logs.} Only the digits
\emph{after} the decimal point in a pH are significant; the part before
it just records the power of ten. $[\ce{H+}] = \num{2.0e-4}$ has two
significant figures, so the pH is written with \textbf{two decimal
places}: 3.70. Writing ``3.7'' understates your precision and ``3.6990''
overstates it.
\end{workedexample}

\begin{yourturn}[YOUR TURN 4 --- four questions]
All at \SI{25}{\celsius}.
\begin{enumerate}[leftmargin=*,itemsep=0.5em]
  \item $[\ce{H+}] = \SI{1.0e-5}{\Molar}$. Find $[\ce{OH-}]$.
        \workspace[1.6cm]
  \item $[\ce{OH-}] = \SI{1.0e-3}{\Molar}$. Find $[\ce{H+}]$ and pH.
        \workspace[1.8cm]
  \item Pure water at \SI{25}{\celsius}: give $[\ce{H+}]$ and pH.
        \workspace[1.6cm]
  \item At \SI{50}{\celsius} neutral water has pH 6.6. Is it acidic?
        Explain. \workspace[1.8cm]
\end{enumerate}
\selfcheck{(a) \SI{1.0e-9}{\Molar} \quad (b) \SI{1.0e-11}{\Molar},
pH 11.00 \quad (c) \SI{1.0e-7}{\Molar}, pH 7.00 \quad (d) no --- it is
neutral}
\keyonly{\begin{apcanswer}(a) $\num{1.0e-14}/\num{1.0e-5} =
\mathbf{\SI{1.0e-9}{\Molar}}$. (b) $[\ce{H+}] =
\num{1.0e-14}/\num{1.0e-3} = \mathbf{\SI{1.0e-11}{\Molar}}$, so
pH $= \mathbf{11.00}$ --- basic, as a large $[\ce{OH-}]$ requires. (c)
$[\ce{H+}] = [\ce{OH-}] = \mathbf{\SI{1.0e-7}{\Molar}}$ and pH
$= \mathbf{7.00}$. (d) \textbf{No --- it is neutral.} Neutral means
$[\ce{H+}] = [\ce{OH-}]$, and that is still true at
\SI{50}{\celsius}. What has changed is $K_{w}$, which grows with
temperature, so the neutral point moves below 7. pH 7 is the neutral
value only at \SI{25}{\celsius}.\end{apcanswer}}
\end{yourturn}

% =====================================================================
\section*{\sffamily Ladder 5 \textbullet\ The pH scale}
\zsec{14.3}

pH is a logarithm, so every whole unit is a \textbf{factor of ten} in
$[\ce{H+}]$.

\begin{center}
\begin{tikzpicture}[font=\sffamily\footnotesize]
  \draw[-{Stealth[length=3mm]},very thick] (0.5,2.3) -- (10.8,2.3);
  \foreach \x/\v in {0.9/0, 2.3/2, 3.7/4, 5.1/6, 5.8/7, 6.5/8,
                     7.9/10, 9.3/12, 10.4/14} {
    \fill (\x,2.3) circle (0.06);
    \node[font=\sffamily\footnotesize] at (\x,2.62) {\v};}
  \node[font=\sffamily\bfseries] at (2.8,1.85) {acidic};
  \node[font=\sffamily\bfseries] at (5.8,1.85) {neutral};
  \node[font=\sffamily\bfseries] at (8.8,1.85) {basic};
  \foreach \x/\lab in {1.3/{stomach acid}, 2.9/{lemon}, 4.6/{coffee},
                       5.8/{pure water}, 7.2/{blood}, 8.6/{ammonia},
                       10.1/{drain cleaner}} {
    \node[rotate=90,anchor=east,font=\sffamily\footnotesize]
      at (\x,1.55) {\lab};}
  \draw[apcgray] (0.4,-0.55) -- (10.8,-0.55);
  \node[anchor=west] at (0.5,-0.95)
    {one pH unit $=$ a \textbf{tenfold} change --- pH 3 is
     \emph{a hundred} times more acidic than pH 5};
\end{tikzpicture}
\end{center}

\begin{workedexample}[I do: what a difference of pH units means]
\textbf{How much more acidic is a pH 2 solution than a pH 5 one?}

The difference is $5 - 2 = 3$ pH units, and each unit is a factor of
ten:
\[ 10^{3} = \mathbf{1000 \text{ times}} \]

\textbf{Verify it directly.} pH 2 means $[\ce{H+}] =
\SI{1e-2}{\Molar}$; pH 5 means $[\ce{H+}] = \SI{1e-5}{\Molar}$. The
ratio is $10^{-2}/10^{-5} = 10^{3}$. Same answer, and this is the
version to fall back on whenever the shortcut feels slippery.

\textbf{The direction that catches people.} \emph{Lower} pH means
\emph{more} acidic, because of the minus sign in $\text{pH} =
-\log[\ce{H+}]$. Higher $[\ce{H+}]$ gives lower pH. If your answer has a
pH 5 solution more acidic than a pH 2 one, you dropped the sign.

\textbf{Going backwards, from pH to concentration.}
\[ [\ce{H+}] = 10^{-\text{pH}} \]
So pH 3.70 gives $[\ce{H+}] = 10^{-3.70} = \SI{2.0e-4}{\Molar}$ ---
which is exactly where Ladder 4's worked example started. The two
operations are inverses, and being fluent in both directions is
assumed from here on.

\textbf{Why the scale runs about 0 to 14.} It is not a hard boundary
--- concentrated strong acids genuinely have negative pH. The range
simply reflects that $K_{w} = 10^{-14}$, so ordinary aqueous solutions
land inside it.
\end{workedexample}

\begin{yourturn}[YOUR TURN 5 --- four questions]
\begin{enumerate}[leftmargin=*,itemsep=0.5em]
  \item How many times more acidic is pH 3 than pH 6?
        \workspace[1.6cm]
  \item A solution has pH 4.00. Find $[\ce{H+}]$.
        \workspace[1.6cm]
  \item A solution has pOH 3.00. Find its pH and say whether it is
        acidic or basic. \workspace[1.8cm]
  \item Which is more acidic, pH 2.5 or pH 4.5? By what factor?
        \workspace[1.8cm]
\end{enumerate}
\selfcheck{(a) 1000 \quad (b) \SI{1.0e-4}{\Molar} \quad (c) pH 11.00,
basic \quad (d) pH 2.5, by 100}
\keyonly{\begin{apcanswer}(a) $10^{3} = \mathbf{1000}$ times.
(b) $[\ce{H+}] = 10^{-4.00} = \mathbf{\SI{1.0e-4}{\Molar}}$.
(c) pH $= 14.00 - 3.00 = \mathbf{11.00}$, which is above 7, so the
solution is \textbf{basic}. (d) \textbf{pH 2.5}, by a factor of
$10^{(4.5-2.5)} = 10^{2} = \mathbf{100}$. Lower pH always means more
acidic.\end{apcanswer}}
\end{yourturn}

% =====================================================================
\section*{\sffamily Ladder 6 \textbullet\ pH of a strong acid}
\zsec{14.4}

The one-line calculation. Complete ionisation means $[\ce{H+}]$ equals
the acid concentration, and you take a logarithm.

\begin{workedexample}[I do: two strong acids]
\textbf{(a) Find the pH of \SI{0.010}{\Molar} \ce{HCl}.}

\ce{HCl} is strong, so ionisation is complete and one \ce{HCl} gives one
\ce{H+}:
\[ [\ce{H+}] = \SI{0.010}{\Molar}
   \qquad
   \text{pH} = -\log(0.010) = \mathbf{2.00} \]

\textbf{(b) Find the pH of \SI{0.0025}{\Molar} \ce{HNO3}.}
\[ [\ce{H+}] = \SI{0.0025}{\Molar}
   \qquad
   \text{pH} = -\log(\num{2.5e-3}) = \mathbf{2.60} \]

\textbf{Check the sig figs.} \num{2.5e-3} has two significant figures,
so the pH carries two decimal places: 2.60, not 2.6 and not 2.602.

\textbf{The trap in this ladder is \ce{H2SO4}.} It is the one strong
acid on the list that is diprotic, and only its \emph{first} proton is
strong --- the second comes off a weak acid, \ce{HSO4-}
($K_{a2} = \num{1.2e-2}$). So \SI{0.010}{\Molar} \ce{H2SO4} gives at
least \SI{0.010}{\Molar} \ce{H+} from the first ionisation, and a bit
more from the second. Unless a question tells you otherwise, treat only
the first proton as complete and say so.

\textbf{And a boundary worth knowing.} For extremely dilute strong acid
--- below about \SI{1e-7}{\Molar} --- water's own autoionisation
contributes as much \ce{H+} as the acid does, and $\text{pH} =
-\log(c)$ breaks down. A \SI{1e-10}{\Molar} \ce{HCl} solution is not
pH 10; it is very slightly below 7. Adding acid can never make water
basic.
\end{workedexample}

\begin{yourturn}[YOUR TURN 6 --- four questions]
\begin{enumerate}[leftmargin=*,itemsep=0.5em]
  \item Find the pH of \SI{0.100}{\Molar} \ce{HCl}.
        \workspace[1.6cm]
  \item Find the pH of \SI{0.0010}{\Molar} \ce{HBr}.
        \workspace[1.6cm]
  \item Find $[\ce{H+}]$ in a strong acid solution of pH 1.50.
        \workspace[1.6cm]
  \item Why can you not simply write $\text{pH} = -\log(c)$ for
        \SI{1e-10}{\Molar} \ce{HCl}? \workspace[1.8cm]
\end{enumerate}
\selfcheck{(a) 1.000 \quad (b) 3.00 \quad (c) \SI{0.032}{\Molar} \quad
(d) water's own \ce{H+} dominates}
\keyonly{\begin{apcanswer}(a) $[\ce{H+}] = \SI{0.100}{\Molar}$, so
pH $= \mathbf{1.000}$ (three sig figs in the concentration, three
decimals in the pH). (b) $-\log(\num{1.0e-3}) = \mathbf{3.00}$.
(c) $[\ce{H+}] = 10^{-1.50} = \mathbf{\SI{0.032}{\Molar}}$. (d) Because
at that concentration the acid supplies far less \ce{H+} than
\textbf{water's own autoionisation} does ($\SI{1e-7}{\Molar}$). The
formula would give pH 10 --- a basic solution produced by adding acid,
which is impossible. The true pH is just below 7.\end{apcanswer}}
\end{yourturn}

% =====================================================================
\section*{\sffamily Ladder 7 \textbullet\ pH of a strong base}
\zsec{14.6}

Same idea, one extra step --- and one place where the formula unit
matters.

\begin{center}
\begin{tikzpicture}[font=\sffamily\footnotesize]
  \node[font=\sffamily\small\bfseries] at (5.3,2.5)
    {count the hydroxides in the formula};
  \node[anchor=west] at (0.5,1.85)
    {$\ce{NaOH -> Na+ + OH-}$ \qquad one \ce{OH-} per formula unit
     \qquad $[\ce{OH-}] = c$};
  \node[anchor=west] at (0.5,1.2)
    {$\ce{Ca(OH)2 -> Ca^2+ + 2OH-}$ \qquad \textbf{two} per formula
     unit \qquad $[\ce{OH-}] = \mathbf{2}c$};
  \draw[apcgray] (0.4,0.8) -- (10.5,0.8);
  \node[anchor=west] at (0.5,0.35)
    {then pOH $= -\log[\ce{OH-}]$, and pH $= 14.00 - $ pOH};
\end{tikzpicture}
\end{center}

\begin{workedexample}[I do: the factor of two]
\textbf{(a) Find the pH of \SI{0.010}{\Molar} \ce{NaOH}.}
\[ [\ce{OH-}] = \SI{0.010}{\Molar}
   \;\Rightarrow\;
   \text{pOH} = 2.00
   \;\Rightarrow\;
   \text{pH} = 14.00 - 2.00 = \mathbf{12.00} \]

\textbf{(b) Find the pH of \SI{0.0050}{\Molar} \ce{Ca(OH)2}.}

Each formula unit releases \emph{two} hydroxides:
\[ [\ce{OH-}] = 2 \times 0.0050 = \SI{0.010}{\Molar} \]
\[ \text{pOH} = 2.00
   \;\Rightarrow\;
   \text{pH} = \mathbf{12.00} \]

\textbf{Notice that (a) and (b) give the same pH} from different
concentrations --- \SI{0.0050}{\Molar} \ce{Ca(OH)2} and
\SI{0.010}{\Molar} \ce{NaOH} are equally basic, because what matters is
$[\ce{OH-}]$, not the concentration written on the bottle. Forgetting
the 2 would have given pOH 2.30 and pH 11.70, an error of 0.30 pH
units.

\textbf{Where the factor does \emph{not} apply.} It applies to strong
bases because they dissociate completely. It does \emph{not} apply to
weak polyprotic acids, whose successive ionisations are separate
equilibria of wildly different strengths (Ladder 13). Complete
dissociation is what lets you simply multiply.

\textbf{One caution about \ce{Ca(OH)2} in practice:} it is only
sparingly soluble, so you cannot actually make a concentrated solution
of it. The calculation is fine for the concentrations exam questions
use.
\end{workedexample}

\begin{yourturn}[YOUR TURN 7 --- four questions]
\begin{enumerate}[leftmargin=*,itemsep=0.5em]
  \item Find the pH of \SI{0.100}{\Molar} \ce{KOH}.
        \workspace[1.6cm]
  \item Find the pH of \SI{0.0010}{\Molar} \ce{Ba(OH)2}.
        \workspace[1.8cm]
  \item Find $[\ce{OH-}]$ in a solution of pH 13.00.
        \workspace[1.6cm]
  \item Why does \ce{Ca(OH)2} need a factor of 2 but \ce{NaOH} does
        not? \workspace[1.6cm]
\end{enumerate}
\selfcheck{(a) 13.000 \quad (b) 11.30 \quad (c) \SI{0.10}{\Molar} \quad
(d) two \ce{OH-} per formula unit}
\keyonly{\begin{apcanswer}(a) $[\ce{OH-}] = \SI{0.100}{\Molar}$,
pOH $= 1.000$, pH $= \mathbf{13.000}$. (b) $[\ce{OH-}] = 2 \times
0.0010 = \SI{0.0020}{\Molar}$; pOH $= -\log(\num{2.0e-3}) = 2.70$; pH
$= 14.00 - 2.70 = \mathbf{11.30}$. (c) pOH $= 14.00 - 13.00 = 1.00$, so
$[\ce{OH-}] = 10^{-1.00} = \mathbf{\SI{0.10}{\Molar}}$. (d) Because one
formula unit of \ce{Ca(OH)2} releases \textbf{two} hydroxide ions on
dissociating, while \ce{NaOH} releases one. The factor counts hydroxides
per formula unit.\end{apcanswer}}
\end{yourturn}

% =====================================================================
\section*{\sffamily Ladder 8 \textbullet\ $K_{a}$ and what ``weak''
measures}
\zsec{14.5}

$K_{a}$ is just an equilibrium constant, applied to one particular
reaction: an acid giving up its proton to water.

\[ \ce{HA <=> H+ + A-}
   \qquad
   K_{a} = \frac{[\ce{H+}][\ce{A-}]}{[\ce{HA}]} \]

\begin{center}
\begin{tikzpicture}[font=\sffamily\footnotesize]
  \draw[-{Stealth[length=3mm]},very thick] (0.6,2.2) -- (10.7,2.2);
  \node[font=\sffamily\bfseries] at (1.6,2.7) {stronger acid};
  \node[font=\sffamily\bfseries] at (9.4,2.7) {weaker acid};
  \foreach \x/\ac/\ka in {%
      1.5/{\ce{HSO4-}}/{$\num{1.2e-2}$},
      3.3/{\ce{HF}}/{$\num{7.2e-4}$},
      5.1/{\ce{CH3COOH}}/{$\num{1.8e-5}$},
      6.9/{\ce{H2CO3}}/{$\num{4.3e-7}$},
      8.7/{\ce{HClO}}/{$\num{3.5e-8}$},
      10.2/{\ce{HCN}}/{$\num{6.2e-10}$}} {
    \fill (\x,2.2) circle (0.07);
    \node[font=\sffamily\footnotesize] at (\x,1.8) {\ac};
    \node[rotate=90,anchor=east,font=\sffamily\footnotesize]
      at (\x,1.45) {\ka};}
  % Rule and caption dropped to -0.5 / -0.95: the rotated Ka values hang
  % down to about y = -0.05, and at 0.15 the rule and caption sat inside
  % them ("6.2e-10" printed over "(Ladder 16)").
  \draw[apcgray] (0.5,-0.5) -- (10.8,-0.5);
  \node[anchor=west] at (0.6,-0.95)
    {\textbf{larger} $K_{a}$ $=$ stronger acid \textbullet\ a
     \textbf{smaller} p$K_{a}$ means the same thing (Ladder 16)};
\end{tikzpicture}
\end{center}

\begin{workedexample}[I do: reading $K_{a}$ values]
\textbf{$K_{a}$ answers ``how far does this acid ionise?'' --- exactly
the question $K$ answered in Chapter 13, for one specific reaction.}

\textbf{Acetic acid, $K_{a} = \num{1.8e-5}$.} Far below 1, so the
equilibrium lies well to the left and the acid is mostly intact in
solution. Weak.

\textbf{\ce{HF}, $K_{a} = \num{7.2e-4}$.} Also below 1, so also weak ---
but about forty times larger than acetic acid's, so \ce{HF} is the
stronger of the two. Compare $K_{a}$ values directly; the bigger number
is the stronger acid, every time.

\textbf{\ce{HCN}, $K_{a} = \num{6.2e-10}$.} Extremely weak. A
\SI{0.1}{\Molar} solution is barely acidic at all.

\textbf{And strong acids?} Their $K_{a}$ values are enormous --- for
\ce{HCl}, roughly $10^{6}$ --- which is why we do not tabulate them.
When $K_{a}$ is that large the equilibrium is so far right that treating
ionisation as complete is not an approximation worth apologising for.

\textbf{The link to the conjugate base, which Ladder 12 makes exact.}
The stronger the acid, the weaker its conjugate base. \ce{HCN} is
feeble, so \ce{CN-} is a genuinely strong base --- strong enough that
sodium cyanide solutions are distinctly alkaline. \ce{HCl} is powerful,
so \ce{Cl-} is so weak a base that it does nothing at all in water.
\end{workedexample}

\begin{yourturn}[YOUR TURN 8 --- four questions]
Use the $K_{a}$ values in the figure above.
\begin{enumerate}[leftmargin=*,itemsep=0.5em]
  \item Which is the stronger acid, \ce{HF} or \ce{HClO}?
        \workspace[1.6cm]
  \item Write the $K_{a}$ expression for \ce{HCN}.
        \workspace[1.6cm]
  \item Which has the stronger conjugate base, \ce{HF} or \ce{HCN}?
        \workspace[1.8cm]
  \item Why are $K_{a}$ values not tabulated for \ce{HCl}?
        \workspace[1.8cm]
\end{enumerate}
\selfcheck{(a) \ce{HF} \quad (b) $[\ce{H+}][\ce{CN-}]/[\ce{HCN}]$ \quad
(c) \ce{HCN} \quad (d) it is strong --- $K_{a}$ is enormous}
\keyonly{\begin{apcanswer}(a) \textbf{\ce{HF}}: $\num{7.2e-4}$ against
\ce{HClO}'s $\num{3.5e-8}$, about twenty thousand times larger. (b)
$K_{a} = \dfrac{[\ce{H+}][\ce{CN-}]}{[\ce{HCN}]}$. (c)
\textbf{\ce{HCN}'s} conjugate base, \ce{CN-}. The weaker the acid, the
stronger its conjugate base --- and \ce{HCN} is by far the weaker of the
two. (d) Because \ce{HCl} is a \textbf{strong} acid: its $K_{a}$ is so
large (around $10^{6}$) that ionisation is effectively complete, and
there is no useful equilibrium to describe. A tabulated value would
never be used.\end{apcanswer}}
\end{yourturn}

% =====================================================================
\section*{\sffamily Ladder 9 \textbullet\ pH of a weak acid}
\zsec{14.5}

The central calculation of Unit 8. It is Chapter 13's ICE table with
$K_{a}$ in place of $K$.

\begin{center}
\begin{tikzpicture}[font=\sffamily\footnotesize]
  \node[font=\sffamily\small\bfseries] at (5.3,3.15)
    {the five steps, every time};
  \foreach \y/\n/\t in {%
      2.55/{1}/{write the ionisation equation and the $K_{a}$
        expression},
      2.05/{2}/{build the ICE table: $c-x$, $x$, $x$},
      1.55/{3}/{assume $x \ll c$, so $K_{a} \approx x^{2}/c$, and solve
        for $x$},
      1.05/{4}/{\textbf{check} $x/c \times 100 < 5\%$},
      0.55/{5}/{pH $= -\log x$, because $x$ \emph{is} $[\ce{H+}]$}} {
    \node[anchor=west,font=\sffamily\bfseries] at (0.5,\y) {\n.};
    \node[anchor=west] at (1.05,\y) {\t};}
  \draw[apcgray] (0.4,2.9) -- (10.5,2.9);
  \draw[apcgray] (0.4,0.2) -- (10.5,0.2);
\end{tikzpicture}
\end{center}

\begin{workedexample}[I do: acetic acid from scratch]
\textbf{Find the pH of \SI{0.500}{\Molar} acetic acid,
$K_{a} = \num{1.8e-5}$.}

\textbf{Step 1 --- equation and expression.}
\[ \ce{CH3COOH <=> H+ + CH3COO-}
   \qquad
   K_{a} = \frac{[\ce{H+}][\ce{CH3COO-}]}{[\ce{CH3COOH}]} \]

\textbf{Step 2 --- ICE table.}
\begin{center}
\renewcommand{\arraystretch}{1.2}
\begin{tabular}{@{}lccc@{}}
\toprule
& $[\ce{CH3COOH}]$ & $[\ce{H+}]$ & $[\ce{CH3COO-}]$ \\
\midrule
\textbf{I} & 0.500 & 0 & 0 \\
\textbf{C} & $-x$ & $+x$ & $+x$ \\
\textbf{E} & $0.500-x$ & $x$ & $x$ \\
\bottomrule
\end{tabular}
\end{center}

\textbf{Step 3 --- substitute and approximate.}
\[ \num{1.8e-5} = \frac{x^{2}}{0.500-x}
   \;\approx\; \frac{x^{2}}{0.500} \]
\[ x^{2} = \num{9.0e-6}
   \;\Longrightarrow\;
   x = \SI{3.0e-3}{\Molar} \]

\textbf{Step 4 --- justify the approximation.}
\[ \frac{\num{3.0e-3}}{0.500} \times 100 = 0.60\%
   \quad<\quad 5\% \quad\checkmark \]

\textbf{Step 5 --- convert to pH.} Because $x$ \emph{is} $[\ce{H+}]$:
\[ \text{pH} = -\log(\num{3.0e-3}) = \mathbf{2.52} \]

\textbf{Sanity check.} A strong acid at \SI{0.500}{\Molar} would have
pH 0.30. Our weak acid comes out at 2.52 --- much less acidic, as it
must be. If a weak-acid pH ever comes out \emph{lower} than the strong
acid equivalent, something is wrong.

\textbf{The single commonest error in this ladder} is stopping at $x$
and reporting it as the answer. $x$ is a concentration; the question
asked for pH. Take the logarithm.
\end{workedexample}

\begin{yourturn}[YOUR TURN 9 --- four questions]
\begin{enumerate}[leftmargin=*,itemsep=0.5em]
  \item Find the pH of \SI{1.00}{\Molar} acetic acid
        ($K_{a} = \num{1.8e-5}$). \workspace[2.0cm]
  \item In the ICE table for a weak acid, what does $x$ represent
        physically? \workspace[1.6cm]
  \item Why must you check the 5\% rule?
        \workspace[1.8cm]
  \item A weak acid and a strong acid are both at
        \SI{0.10}{\Molar}. Which has the lower pH, and why?
        \workspace[1.8cm]
\end{enumerate}
\selfcheck{(a) 2.37 \quad (b) $[\ce{H+}]$ at equilibrium \quad (c) to
confirm dropping $x$ was legitimate \quad (d) the strong acid}
\keyonly{\begin{apcanswer}(a) $x^{2}/1.00 \approx \num{1.8e-5}$, so
$x = \num{4.2e-3}$ and pH $= -\log(\num{4.2e-3}) = \mathbf{2.37}$.
(Check: $\num{4.2e-3}/1.00 = 0.42\% < 5\%$ \checkmark) (b) $x$ is the
\textbf{equilibrium $[\ce{H+}]$} --- and also the equilibrium
$[\ce{A-}]$, since the acid produces them one for one. (c) Because
dropping $x$ from $c - x$ is only valid when $x$ is genuinely negligible
against $c$. If the check fails, the simplification introduced real
error and you must solve the quadratic instead. Examiners award the mark
for performing the check. (d) The \textbf{strong acid}. It ionises
completely, giving $[\ce{H+}] = \SI{0.10}{\Molar}$ and pH 1.00, while
the weak acid gives far less \ce{H+} and therefore a higher
pH.\end{apcanswer}}
\end{yourturn}

% =====================================================================
\section*{\sffamily Ladder 10 \textbullet\ Percent ionisation}
\zsec{14.5}

\[ \text{percent ionisation}
   = \frac{[\ce{H+}]_{\text{equilibrium}}}{[\ce{HA}]_{0}} \times 100 \]

This is the same fraction as the 5\% check --- but here it is the
answer, not a validity test.

\begin{center}
\begin{tikzpicture}
\begin{axis}[
  width=10.4cm, height=5.6cm,
  xlabel={$[\ce{CH3COOH}]_{0}$ (\si{\Molar})},
  ylabel={percent ionisation},
  xmin=0, xmax=1.05, ymin=0, ymax=5.2,
  xtick={0,0.2,0.4,0.6,0.8,1.0}, ytick={0,1,2,3,4,5},
  grid=major, grid style={apcgray!25},
  label style={font=\footnotesize}, tick label style={font=\footnotesize},
]
\addplot[seriesA, thick, domain=0.01:1.0, samples=90]
  {100*sqrt(1.8e-5/x)};
\addplot[only marks, mark=*, mark size=1.6pt, black]
  coordinates {(0.010,4.24) (0.100,1.34) (1.00,0.42)};
\node[font=\footnotesize,anchor=west] at (axis cs:0.09,4.55)
  {more dilute $\Rightarrow$ \textbf{more} ionised};
\node[font=\footnotesize,anchor=west] at (axis cs:0.30,2.6)
  {same acid throughout --- only the};
\node[font=\footnotesize,anchor=west] at (axis cs:0.30,2.15)
  {concentration changes};
\end{axis}
\end{tikzpicture}
\end{center}

\begin{apctrap}
\textbf{Diluting a weak acid raises its percent ionisation but
\emph{lowers} its $[\ce{H+}]$.} Both are true at once and they sound
contradictory. A larger \emph{fraction} of a smaller \emph{amount} still
gives less acid in total --- so the solution gets less acidic (pH rises)
even as the percentage climbs. And $K_{a}$ itself never moves; only
temperature changes that.
\end{apctrap}

\begin{workedexample}[I do: the same acid at three concentrations]
\textbf{Acetic acid, $K_{a} = \num{1.8e-5}$, at three
concentrations.}

\textbf{\SI{1.00}{\Molar}:} $x = \sqrt{\num{1.8e-5} \times 1.00} =
\num{4.24e-3}$, so
\[ \frac{\num{4.24e-3}}{1.00} \times 100 = \mathbf{0.42\%} \]

\textbf{\SI{0.100}{\Molar}:} $x = \num{1.34e-3}$, so
\[ \frac{\num{1.34e-3}}{0.100} \times 100 = \mathbf{1.3\%} \]

\textbf{\SI{0.0100}{\Molar}:} $x = \num{4.24e-4}$, so
\[ \frac{\num{4.24e-4}}{0.0100} \times 100 = \mathbf{4.2\%} \]

\textbf{Read the pattern.} A tenfold dilution multiplies the percent
ionisation by about $\sqrt{10} \approx 3.2$ --- from 0.42 to 1.3 to 4.2.
That square-root behaviour follows straight from $x = \sqrt{K_{a}c}$:
percent ionisation is $x/c = \sqrt{K_{a}/c}$, so it rises as $c$ falls.

\textbf{But check what happened to the acidity.} $[\ce{H+}]$ went
\emph{down} at every step, from \num{4.24e-3} to \num{4.24e-4}, so the
pH rose from about 2.4 to about 3.4. The solution became \emph{less}
acidic while becoming \emph{more} ionised. Both statements are correct,
and a question that asks for one while implying the other is testing
whether you can hold them apart.

\textbf{Why dilution pushes ionisation up.} Le Ch\^{a}telier, from
Chapter 13: adding water dilutes everything, and the equilibrium shifts
toward the side with \emph{more} dissolved particles --- the ionised
side, since one \ce{HA} becomes two ions.
\end{workedexample}

\begin{yourturn}[YOUR TURN 10 --- four questions]
\begin{enumerate}[leftmargin=*,itemsep=0.5em]
  \item A \SI{0.20}{\Molar} weak acid has $[\ce{H+}] =
        \SI{2.0e-3}{\Molar}$. Find its percent ionisation.
        \workspace[1.8cm]
  \item Does diluting a weak acid raise or lower its percent
        ionisation? \workspace[1.6cm]
  \item Does diluting a weak acid raise or lower its pH?
        \workspace[1.6cm]
  \item Does diluting a weak acid change its $K_{a}$? Explain.
        \workspace[1.8cm]
\end{enumerate}
\selfcheck{(a) 1.0\% \quad (b) raises it \quad (c) raises it (less
acidic) \quad (d) no --- only temperature does}
\keyonly{\begin{apcanswer}(a) $\dfrac{\num{2.0e-3}}{0.20} \times 100 =
\mathbf{1.0\%}$. (b) \textbf{Raises} it --- a larger fraction of the
acid comes apart. (c) \textbf{Raises} the pH, meaning the solution
becomes \emph{less} acidic. $[\ce{H+}]$ falls even though the
\emph{percentage} ionised rises, because a bigger fraction of a much
smaller amount is still less acid. (d) \textbf{No.} $K_{a}$ is an
equilibrium constant and depends only on \textbf{temperature}. Dilution
moves the position of the equilibrium, not the constant that describes
it.\end{apcanswer}}
\end{yourturn}

% =====================================================================
\section*{\sffamily Ladder 11 \textbullet\ Weak bases and $K_{b}$}
\zsec{14.6}

Identical machinery, one substitution: the base takes a proton
\emph{from water}, so what you solve for is $[\ce{OH-}]$, and you must
finish by converting to pH.

\[ \ce{B + H2O <=> BH+ + OH-}
   \qquad
   K_{b} = \frac{[\ce{BH+}][\ce{OH-}]}{[\ce{B}]} \]

\begin{workedexample}[I do: ammonia, and the extra step]
\textbf{Find the pH of \SI{0.500}{\Molar} \ce{NH3},
$K_{b} = \num{1.8e-5}$.}

\textbf{ICE gives the same algebra as Ladder 9:}
\[ \num{1.8e-5} \approx \frac{x^{2}}{0.500}
   \;\Longrightarrow\;
   x = \SI{3.0e-3}{\Molar} \]
Check: $0.60\% < 5\%$ \checkmark

\textbf{Now the step that separates bases from acids.} Here $x$ is
$[\ce{OH-}]$, not $[\ce{H+}]$, so it gives you \textbf{pOH} first:
\[ \text{pOH} = -\log(\num{3.0e-3}) = 2.52 \]
\[ \text{pH} = 14.00 - 2.52 = \mathbf{11.48} \]

\textbf{Compare with Ladder 9 deliberately.} Acetic acid at
\SI{0.500}{\Molar} with $K_{a} = \num{1.8e-5}$ gave pH 2.52. Ammonia at
\SI{0.500}{\Molar} with the same numerical $K_{b}$ gives pH 11.48. The
two are mirror images about 7 --- $2.52 + 11.48 = 14.00$ --- because the
constants are numerically equal. That symmetry is worth noticing, and it
is a fast way to spot a forgotten conversion.

\textbf{Which is exactly the error to watch for.} Stopping at
$\text{pH} = -\log x = 2.52$ would report a strongly \emph{acidic}
ammonia solution. Ammonia is a base; its pH must exceed 7. Always ask
whether the sign of your answer matches the chemistry before you write
it down.
\end{workedexample}

\begin{yourturn}[YOUR TURN 11 --- four questions]
\begin{enumerate}[leftmargin=*,itemsep=0.5em]
  \item Write the $K_{b}$ expression for methylamine, \ce{CH3NH2}.
        \workspace[1.8cm]
  \item For \SI{0.100}{\Molar} of a base with $K_{b} = \num{1.0e-5}$,
        find $[\ce{OH-}]$. \workspace[1.8cm]
  \item Convert that $[\ce{OH-}]$ to a pH. \workspace[1.8cm]
  \item Why does a weak-base ICE table give pOH rather than pH
        directly? \workspace[1.8cm]
\end{enumerate}
\selfcheck{(a) $[\ce{CH3NH3+}][\ce{OH-}]/[\ce{CH3NH2}]$ \quad (b)
\SI{1.0e-3}{\Molar} \quad (c) 11.00 \quad (d) $x$ is $[\ce{OH-}]$}
\keyonly{\begin{apcanswer}(a) $K_{b} =
\dfrac{[\ce{CH3NH3+}][\ce{OH-}]}{[\ce{CH3NH2}]}$. (b)
$x^{2}/0.100 \approx \num{1.0e-5}$, so $x^{2} = \num{1.0e-6}$ and
$[\ce{OH-}] = \mathbf{\SI{1.0e-3}{\Molar}}$. (c) pOH $=
-\log(\num{1.0e-3}) = 3.00$, so pH $= 14.00 - 3.00 = \mathbf{11.00}$ ---
basic, as required. (d) Because the base reacts with \textbf{water},
producing \ce{OH-}. The $x$ you solve for is therefore
\textbf{$[\ce{OH-}]$}, which gives pOH directly; pH requires the extra
step pH $= 14.00 - \text{pOH}$.\end{apcanswer}}
\end{yourturn}

% =====================================================================
\section*{\sffamily Ladder 12 \textbullet\ $K_{a} \times K_{b} = K_{w}$}
\zsec{14.6}

For a \textbf{conjugate pair} --- and only for a conjugate pair --- the
two constants multiply to $K_{w}$.

\[ K_{a} \times K_{b} = K_{w} = \num{1.0e-14}
   \qquad\text{equivalently}\qquad
   \text{p}K_{a} + \text{p}K_{b} = 14.00 \]

\begin{center}
\begin{tikzpicture}[font=\sffamily\footnotesize]
  \node[font=\sffamily\small\bfseries] at (5.3,2.75)
    {stronger acid $\Longleftrightarrow$ weaker conjugate base};
  \foreach \y/\ac/\ka/\cb/\kb in {%
      2.1/{\ce{HF}}/{$\num{7.2e-4}$}/{\ce{F-}}/{$\num{1.4e-11}$},
      1.5/{\ce{CH3COOH}}/{$\num{1.8e-5}$}/{\ce{CH3COO-}}/{$\num{5.6e-10}$},
      0.9/{\ce{HCN}}/{$\num{6.2e-10}$}/{\ce{CN-}}/{$\num{1.6e-5}$}} {
    \node[anchor=west] at (0.5,\y) {\ac};
    \node[anchor=west] at (2.6,\y) {$K_{a} = $ \ka};
    \node[anchor=west] at (5.5,\y) {\cb};
    \node[anchor=west] at (7.6,\y) {$K_{b} = $ \kb};}
  \draw[apcgray] (0.4,2.5) -- (10.6,2.5);
  \draw[apcgray] (0.4,0.55) -- (10.6,0.55);
  \node[anchor=west] at (0.5,0.1)
    {each row multiplies to $\num{1.0e-14}$ --- as $K_{a}$ falls down
     the table, $K_{b}$ rises};
\end{tikzpicture}
\end{center}

\begin{workedexample}[I do: converting between the two]
\textbf{Acetic acid has $K_{a} = \num{1.8e-5}$. Find $K_{b}$ for
acetate, \ce{CH3COO-}.}
\[ K_{b} = \frac{K_{w}}{K_{a}}
         = \frac{\num{1.0e-14}}{\num{1.8e-5}}
         = \mathbf{\num{5.6e-10}} \]

\textbf{Sense-check the size.} $\num{5.6e-10}$ is small, so acetate is a
weak base --- but not a negligible one, and that is precisely why sodium
acetate solutions come out basic (Ladder 14).

\textbf{Now the other direction. Ammonia has $K_{b} = \num{1.8e-5}$;
find $K_{a}$ for the ammonium ion, \ce{NH4+}.}
\[ K_{a} = \frac{\num{1.0e-14}}{\num{1.8e-5}} = \num{5.6e-10} \]
So \ce{NH4+} is a weak acid --- which is why ammonium chloride
solutions are acidic.

\textbf{The relationship must be used on a conjugate pair.} $K_{a}$ for
acetic acid and $K_{b}$ for ammonia do \emph{not} multiply to $K_{w}$;
they are unrelated species. Pair the acid with \emph{its own} conjugate
base, differing by exactly one proton.

\textbf{Why the rule holds at all.} Add the two equilibria:
\[ \ce{HA <=> H+ + A-}
   \qquad
   \ce{A- + H2O <=> HA + OH-} \]
Everything cancels except $\ce{H2O <=> H+ + OH-}$, whose constant is
$K_{w}$ --- and from Chapter 13, adding equations multiplies their
constants. It is Ladder 7 of Chapter 13, reused.
\end{workedexample}

\begin{yourturn}[YOUR TURN 12 --- four questions]
\begin{enumerate}[leftmargin=*,itemsep=0.5em]
  \item \ce{HF} has $K_{a} = \num{7.2e-4}$. Find $K_{b}$ for \ce{F-}.
        \workspace[1.8cm]
  \item A base has $K_{b} = \num{1.0e-6}$. Find $K_{a}$ for its
        conjugate acid. \workspace[1.8cm]
  \item An acid has p$K_{a} = 4.00$. Find p$K_{b}$ for its conjugate
        base. \workspace[1.6cm]
  \item Why do $K_{a}$ for acetic acid and $K_{b}$ for ammonia
        \emph{not} multiply to $K_{w}$? \workspace[1.8cm]
\end{enumerate}
\selfcheck{(a) \num{1.4e-11} \quad (b) \num{1.0e-8} \quad (c) 10.00
\quad (d) they are not a conjugate pair}
\keyonly{\begin{apcanswer}(a) $K_{b} = K_{w}/K_{a}$, which is
$\num{1.0e-14}/\num{7.2e-4} = \mathbf{\num{1.4e-11}}$. (b) The same
move gives $\num{1.0e-14}/\num{1.0e-6} = \mathbf{\num{1.0e-8}}$.
(c) $14.00 - 4.00 = \mathbf{10.00}$. (d) Because
they are \textbf{not a conjugate pair} --- acetate is the conjugate base
of acetic acid, and ammonia is not. The relationship only holds for two
species differing by exactly one proton.\end{apcanswer}}
\end{yourturn}

% =====================================================================
\section*{\sffamily Ladder 13 \textbullet\ Polyprotic acids}
\zsec{14.7}

Acids with more than one ionisable proton give them up in stages, and
each stage is far weaker than the one before.

\begin{center}
\begin{tikzpicture}[font=\sffamily\footnotesize]
  \node[font=\sffamily\small\bfseries] at (5.3,3.0)
    {phosphoric acid: each step is $\sim$$10^{5}$ times weaker};
  \foreach \y/\eq/\ka in {%
      2.35/{$\ce{H3PO4 <=> H+ + H2PO4^-}$}/{$K_{a1} = \num{7.5e-3}$},
      1.75/{$\ce{H2PO4^- <=> H+ + HPO4^2-}$}/{$K_{a2} = \num{6.2e-8}$},
      1.15/{$\ce{HPO4^2- <=> H+ + PO4^3-}$}/{$K_{a3} = \num{4.8e-13}$}} {
    \node[anchor=west] at (0.6,\y) {\eq};
    \node[anchor=west] at (6.2,\y) {\ka};}
  \draw[apcgray] (0.4,2.75) -- (10.5,2.75);
  \draw[apcgray] (0.4,0.8) -- (10.5,0.8);
  \node[anchor=west] at (0.5,0.35)
    {so the \textbf{first} ionisation supplies essentially all the
     \ce{H+} --- use $K_{a1}$ alone for the pH};
\end{tikzpicture}
\end{center}

\begin{workedexample}[I do: why only the first step matters]
\textbf{Why does $K_{a2}$ drop so far below $K_{a1}$?}

\textbf{Because of charge.} Pulling \ce{H+} away from neutral
\ce{H3PO4} is one thing. Pulling a \emph{positive} \ce{H+} away from an
already \emph{negative} \ce{H2PO4-} means fighting an electrostatic
attraction --- and it gets harder again for the $2-$ ion. Each
successive proton leaves behind a more negative species, so each is
harder to remove.

\textbf{The consequence for calculations.} $K_{a1}$ exceeds $K_{a2}$ by
about $10^{5}$, so the second ionisation contributes a hundred-thousandth
as much \ce{H+} as the first. For pH purposes you can ignore everything
after the first step and treat the acid as an ordinary weak monoprotic
acid.

\textbf{Worked: the pH of \SI{0.100}{\Molar} \ce{H2CO3},
$K_{a1} = \num{4.3e-7}$.}
\[ x^{2}/0.100 \approx \num{4.3e-7}
   \;\Longrightarrow\;
   x^{2} = \num{4.3e-8}
   \;\Longrightarrow\;
   x = \num{2.07e-4} \]
\[ \text{pH} = -\log(\num{2.07e-4}) = \mathbf{3.68} \]
Check: $\num{2.07e-4}/0.100 = 0.21\% < 5\%$ \checkmark. And
$K_{a2} = \num{5.6e-11}$ is nearly four orders of magnitude smaller, so
ignoring it was safe.

\textbf{The one exception on the strong-acid list.} \ce{H2SO4}'s first
ionisation is \emph{strong} and its second is weak but not negligible
($K_{a2} = \num{1.2e-2}$). It is the polyprotic acid where the second
step genuinely matters, and it is the one an exam question will use if
it wants to test this.
\end{workedexample}

\begin{yourturn}[YOUR TURN 13 --- four questions]
\begin{enumerate}[leftmargin=*,itemsep=0.5em]
  \item Why is $K_{a2}$ always smaller than $K_{a1}$?
        \workspace[1.8cm]
  \item Which $K$ do you use to find the pH of \SI{0.10}{\Molar}
        \ce{H3PO4}? \workspace[1.6cm]
  \item Write the equation for the second ionisation of \ce{H2CO3}.
        \workspace[1.6cm]
  \item Why is \ce{H2SO4} the exception among polyprotic acids?
        \workspace[1.8cm]
\end{enumerate}
\selfcheck{(a) removing \ce{H+} from a negative ion is harder \quad (b)
$K_{a1}$ \quad (c) $\ce{HCO3^- <=> H+ + CO3^2-}$ \quad (d) its first
ionisation is strong}
\keyonly{\begin{apcanswer}(a) Because the second proton must be pulled
away from an \textbf{already negatively charged} species, and the
electrostatic attraction between the \ce{H+} and that negative ion makes
removal much harder. Each successive step leaves a more negative
residue. (b) \textbf{$K_{a1}$ alone.} It is roughly $10^{5}$ times
larger than $K_{a2}$, so the first ionisation supplies essentially all
the \ce{H+}. (c) $\ce{HCO3^- <=> H+ + CO3^2-}$. (d) Because its
\textbf{first ionisation is strong} (complete), while its second is a
weak equilibrium with $K_{a2} = \num{1.2e-2}$ --- large enough that it
cannot simply be ignored. Every other common polyprotic acid is weak at
every step.\end{apcanswer}}
\end{yourturn}

% =====================================================================
\section*{\sffamily Ladder 14 \textbullet\ Acid--base properties of
salts}
\zsec{14.8}

Dissolve a salt and its ions may react with water. Whether the solution
comes out acidic, basic or neutral depends on which parent acid and base
the salt came from.

\begin{center}
\begin{tikzpicture}[font=\sffamily\footnotesize]
  \node[font=\sffamily\small\bfseries] at (5.4,3.6)
    {trace each ion back to its parent};
  \node[font=\sffamily\bfseries,align=center] at (3.5,2.95)
    {from a \textbf{strong} acid};
  \node[font=\sffamily\bfseries,align=center] at (7.2,2.95)
    {from a \textbf{weak} acid};
  \node[anchor=east,font=\sffamily\bfseries,align=right] at (2.0,2.15)
    {from a \textbf{strong} base};
  \node[anchor=east,font=\sffamily\bfseries,align=right] at (2.0,1.15)
    {from a \textbf{weak} base};
  \draw[apcgray] (2.2,2.6) -- (9.2,2.6);
  \draw[apcgray] (2.2,1.65) -- (9.2,1.65);
  \draw[apcgray] (5.35,2.75) -- (5.35,0.6);
  \node[align=center] at (3.5,2.15) {\textbf{neutral}\\\ce{NaCl}};
  \node[align=center] at (7.2,2.15) {\textbf{basic}\\\ce{NaCH3COO}};
  \node[align=center] at (3.5,1.15) {\textbf{acidic}\\\ce{NH4Cl}};
  \node[align=center] at (7.2,1.15) {compare $K_{a}$, $K_{b}$
    \\\ce{NH4CH3COO}};
  \node[font=\sffamily\footnotesize] at (5.4,0.15)
    {the ion from the \emph{weak} parent is the one that reacts ---
     the other is a spectator};
\end{tikzpicture}
\end{center}

\begin{workedexample}[I do: three salts, one method]
\textbf{The method: split the salt into its ions, ask where each came
from, and let the ion from the weak parent decide.}

\textbf{\ce{NaCl}.} \ce{Na+} is from \ce{NaOH}, a strong base;
\ce{Cl-} is from \ce{HCl}, a strong acid. Both conjugates are so weak
they do nothing at all. \textbf{Neutral}, pH 7.

\textbf{\ce{NH4Cl}.} \ce{Cl-} is again a spectator. \ce{NH4+} is the
conjugate acid of the \emph{weak} base ammonia, so it is a weak acid and
donates a proton to water. \textbf{Acidic}.

\textbf{\ce{NaCH3COO}.} \ce{Na+} is a spectator. Acetate is the
conjugate base of the \emph{weak} acid acetic acid, so it takes a
proton from water and releases \ce{OH-}. \textbf{Basic}.

\textbf{Worked calculation.} Find the pH of a \SI{0.10}{\Molar}
solution of sodium acetate. From Ladder 12 we have
$K_{b}(\ce{CH3COO-}) = \num{5.6e-10}$, so:
\[ x^{2}/0.10 \approx \num{5.6e-10}
   \;\Longrightarrow\;
   x = \num{7.5e-6} = [\ce{OH-}] \]
\[ \text{pOH} = 5.13
   \;\Longrightarrow\;
   \text{pH} = 14.00 - 5.13 = \mathbf{8.87} \]
Basic, as predicted --- and only mildly so, which is what a $K_{b}$ of
$10^{-10}$ should give.

\textbf{The fourth box.} When \emph{both} ions come from weak parents,
as in ammonium acetate, you must compare $K_{a}$ of the cation with
$K_{b}$ of the anion; the larger one wins. For ammonium acetate both are
$\num{5.6e-10}$, so the solution is very nearly neutral --- a
coincidence of that particular pair, not a general rule.

\textbf{One more class worth recognising:} small, highly charged metal
cations such as \ce{Al^3+} and \ce{Fe^3+} make solutions acidic, because
they polarise the water molecules bound to them enough to release
\ce{H+}. \ce{AlCl3} solutions are noticeably acidic despite containing
no obvious acid.
\end{workedexample}

\begin{yourturn}[YOUR TURN 14 --- four questions]
\begin{enumerate}[leftmargin=*,itemsep=0.5em]
  \item Is a \ce{KNO3} solution acidic, basic or neutral? Explain.
        \workspace[1.8cm]
  \item Is a \ce{NaCN} solution acidic, basic or neutral? Explain.
        \workspace[1.8cm]
  \item Is a \ce{NH4NO3} solution acidic, basic or neutral? Explain.
        \workspace[1.8cm]
  \item Which ion in \ce{NaF} reacts with water, and what does it
        produce? \workspace[1.8cm]
\end{enumerate}
\selfcheck{(a) neutral \quad (b) basic \quad (c) acidic \quad (d)
\ce{F-}, producing \ce{OH-}}
\keyonly{\begin{apcanswer}(a) \textbf{Neutral.} \ce{K+} comes from the
strong base \ce{KOH} and \ce{NO3-} from the strong acid \ce{HNO3}; both
are spectators. (b) \textbf{Basic.} \ce{CN-} is the conjugate base of
the very weak acid \ce{HCN}, so it is a relatively strong base and takes
a proton from water, releasing \ce{OH-}. (c) \textbf{Acidic.} \ce{NH4+}
is the conjugate acid of the weak base ammonia and donates a proton to
water; \ce{NO3-} is a spectator. (d) \textbf{\ce{F-}}, the conjugate
base of the weak acid \ce{HF}. It reacts as
$\ce{F- + H2O <=> HF + OH-}$, producing \textbf{\ce{OH-}} and making
the solution basic.\end{apcanswer}}
\end{yourturn}

% =====================================================================
\section*{\sffamily Ladder 15 \textbullet\ Structure and acid strength}
\zsec{14.9--14.10}

Why is \ce{HI} strong and \ce{HF} weak? Why is \ce{HClO4} the strongest
common acid? Two families, two different explanations.

\begin{center}
\begin{tikzpicture}[font=\sffamily\footnotesize]
  \node[font=\sffamily\small\bfseries] at (5.3,3.5)
    {binary acids \ce{H-X}};
  \node[anchor=west] at (0.5,2.9)
    {\textbf{down} a group: \quad $\ce{HF} < \ce{HCl} < \ce{HBr} <
     \ce{HI}$ \quad --- the \ce{H-X} \textbf{bond gets weaker}};
  \node[anchor=west] at (0.5,2.35)
    {\textbf{across} a period: \quad $\ce{CH4} < \ce{NH3} < \ce{H2O} <
     \ce{HF}$ \quad --- X gets \textbf{more electronegative}};
  \draw[apcgray] (0.4,1.95) -- (10.5,1.95);
  \node[font=\sffamily\small\bfseries] at (5.3,1.55)
    {oxyacids};
  \node[anchor=west] at (0.5,1.0)
    {\textbf{more oxygens}: \quad $\ce{HClO} < \ce{HClO2} < \ce{HClO3}
     < \ce{HClO4}$};
  \node[anchor=west] at (0.5,0.45)
    {\textbf{same oxygens, more electronegative centre}: \quad
     $\ce{HIO} < \ce{HBrO} < \ce{HClO}$};
  \draw[apcgray] (0.4,0.1) -- (10.5,0.1);
  \node[anchor=west] at (0.5,-0.35)
    {both oxyacid rules work the same way: pull electron density
     \emph{away} from the \ce{O-H} bond};
\end{tikzpicture}
\end{center}

\begin{workedexample}[I do: two rules, two reasons]
\textbf{Binary acids down a group: \ce{HF} $<$ \ce{HCl} $<$ \ce{HBr}
$<$ \ce{HI}.}

You might expect fluorine, the most electronegative element, to give the
strongest acid. It does not, and the reason is \textbf{bond strength}.
Iodine is large, so the \ce{H-I} bond is long and weak and the proton
leaves easily. The \ce{H-F} bond is short and very strong, and holding
onto the proton is exactly what makes \ce{HF} weak. Going down a group,
bond strength beats electronegativity.

\textbf{Binary acids across a period: \ce{CH4} $<$ \ce{NH3} $<$
\ce{H2O} $<$ \ce{HF}.} Here the sizes are similar, so bond strength no
longer dominates and \textbf{electronegativity} takes over. A more
electronegative X pulls the bonding electrons toward itself, leaving the
hydrogen more positive and easier to remove --- and stabilising the
anion left behind.

\textbf{Oxyacids, more oxygens: \ce{HClO} $<$ \ce{HClO2} $<$ \ce{HClO3}
$<$ \ce{HClO4}.} Each extra oxygen is highly electronegative and pulls
electron density away from the \ce{O-H} bond, weakening it. It also
spreads the negative charge of the anion over more atoms, which
stabilises the conjugate base. Both effects push in the same direction.

\textbf{Oxyacids, same oxygens: \ce{HIO} $<$ \ce{HBrO} $<$ \ce{HClO}.}
Same argument with the central atom instead: chlorine is the most
electronegative of the three and withdraws the most.

\textbf{Say it in terms of the conjugate base.} ``The more the negative
charge is delocalised, the more stable the conjugate base, so the
stronger the acid.'' That single sentence answers most structure
questions in this ladder, and it earns marks where ``it has more
oxygens'' does not.
\end{workedexample}

\begin{yourturn}[YOUR TURN 15 --- four questions]
\begin{enumerate}[leftmargin=*,itemsep=0.5em]
  \item Which is the stronger acid, \ce{HBr} or \ce{HCl}? Why?
        \workspace[1.8cm]
  \item Which is the stronger acid, \ce{HClO3} or \ce{HClO}? Why?
        \workspace[1.8cm]
  \item Which is the stronger acid, \ce{H2O} or \ce{NH3}? Why?
        \workspace[1.8cm]
  \item Why is \ce{HF} weak when \ce{F} is the most electronegative
        element? \workspace[1.8cm]
\end{enumerate}
\selfcheck{(a) \ce{HBr} --- weaker bond \quad (b) \ce{HClO3} --- more
oxygens \quad (c) \ce{H2O} --- O is more electronegative \quad (d) the
\ce{H-F} bond is very strong}
\keyonly{\begin{apcanswer}(a) \textbf{\ce{HBr}}. Bromine is larger than
chlorine, so the \ce{H-Br} bond is longer and \textbf{weaker} and the
proton is released more easily. Down a group, bond strength decides.
(b) \textbf{\ce{HClO3}}. Its two extra oxygens withdraw electron density
from the \ce{O-H} bond and \textbf{delocalise the negative charge} of
the conjugate base over more atoms, stabilising it. (c)
\textbf{\ce{H2O}}. Across a period the atoms are of similar size, so
\textbf{electronegativity} decides, and oxygen is more electronegative
than nitrogen. (d) Because down a group \textbf{bond strength matters
more than electronegativity}. The \ce{H-F} bond is short and unusually
strong, so despite fluorine's electronegativity the proton is held
tightly and \ce{HF} does not ionise completely.\end{apcanswer}}
\end{yourturn}

% =====================================================================
\section*{\sffamily Ladder 16 \textbullet\ pH and p$K_{a}$}
\zsec{14.8}

p$K_{a} = -\log K_{a}$, exactly as pH is $-\log[\ce{H+}]$. Comparing the
two tells you which form of the acid dominates.

\begin{center}
\begin{tikzpicture}[font=\sffamily\footnotesize]
  \draw[-{Stealth[length=3mm]},very thick] (0.6,2.3) -- (10.6,2.3);
  \node[font=\sffamily\bfseries] at (5.6,2.75) {p$K_{a}$};
  \fill (5.6,2.3) circle (0.11);
  \draw[densely dashed,apcgray] (5.6,2.02) -- (5.6,2.25);
  \node[font=\sffamily\bfseries] at (2.6,1.85) {pH $<$ p$K_{a}$};
  \node[align=center] at (2.6,1.15)
    {more acidic than the acid\\\textbf{\ce{HA}} dominates};
  \node[font=\sffamily\bfseries] at (5.6,1.85) {pH $=$ p$K_{a}$};
  \node[align=center] at (5.6,1.15) {equal amounts\\of each};
  \node[font=\sffamily\bfseries] at (8.7,1.85) {pH $>$ p$K_{a}$};
  \node[align=center] at (8.7,1.15)
    {less acidic than the acid\\\textbf{\ce{A-}} dominates};
  \draw[apcgray] (0.5,0.6) -- (10.7,0.6);
  \node[anchor=west] at (0.6,0.15)
    {\textbf{lower} p$K_{a}$ $=$ \textbf{stronger} acid --- the minus
     sign flips the order, exactly as it does for pH};
\end{tikzpicture}
\end{center}

\begin{workedexample}[I do: p$K_{a}$ and which form you have]
\textbf{Acetic acid has $K_{a} = \num{1.8e-5}$. Find p$K_{a}$.}
\[ \text{p}K_{a} = -\log(\num{1.8e-5}) = \mathbf{4.74} \]

\textbf{Now compare acids by p$K_{a}$.} \ce{HF} has
$K_{a} = \num{7.2e-4}$, so p$K_{a} = 3.14$; \ce{HCN} has
$K_{a} = \num{6.2e-10}$, so p$K_{a} = 9.21$. The order of strength is
$\ce{HF} > \ce{CH3COOH} > \ce{HCN}$, and the p$K_{a}$ values 3.14,
4.74, 9.21 rise in exactly the reverse order. \textbf{Lower p$K_{a}$,
stronger acid.}

\textbf{Now use pH against p$K_{a}$.} Put acetic acid
(p$K_{a} = 4.74$) into a solution buffered at pH 2.00. Since pH $<$
p$K_{a}$, the solution is more acidic than the acid's own crossover
point, and the \textbf{protonated} form \ce{CH3COOH} dominates.

At pH 7.00, pH $>$ p$K_{a}$, so the \textbf{deprotonated} form
\ce{CH3COO-} dominates. And at pH 4.74 exactly, the two are present in
equal amounts.

\textbf{Why this matters beyond the exam.} It is how a drug's behaviour
is predicted: aspirin has p$K_{a} \approx 3.5$, so in stomach acid
(pH $\approx 2$) it is neutral and crosses membranes easily, while in
blood (pH 7.4) it is ionised and stays put. Same molecule, two
environments, opposite behaviour --- read off a comparison of pH with
p$K_{a}$.

\textbf{The rule in one line:} \emph{low pH protonates, high pH
deprotonates}, and p$K_{a}$ is the pH at which the balance tips.
\end{workedexample}

\begin{yourturn}[YOUR TURN 16 --- four questions]
\begin{enumerate}[leftmargin=*,itemsep=0.5em]
  \item An acid has $K_{a} = \num{1.0e-3}$. Find p$K_{a}$.
        \workspace[1.6cm]
  \item Acid A has p$K_{a} = 3.0$ and acid B has p$K_{a} = 6.0$. Which
        is stronger? \workspace[1.6cm]
  \item An acid with p$K_{a} = 5.0$ sits in a solution at pH 8.0. Which
        form dominates? \workspace[1.8cm]
  \item At what pH are $[\ce{HA}]$ and $[\ce{A-}]$ equal?
        \workspace[1.6cm]
\end{enumerate}
\selfcheck{(a) 3.00 \quad (b) A \quad (c) \ce{A-} \quad (d) pH $=$
p$K_{a}$}
\keyonly{\begin{apcanswer}(a) $-\log(\num{1.0e-3}) = \mathbf{3.00}$.
(b) \textbf{A.} A lower p$K_{a}$ means a larger $K_{a}$ and therefore a
stronger acid. (c) \textbf{\ce{A-}}, the deprotonated form. pH 8.0 is
above the p$K_{a}$ of 5.0, so the solution is less acidic than the
crossover point and the conjugate base dominates. (d) When
\textbf{pH $=$ p$K_{a}$} --- that is what p$K_{a}$
marks.\end{apcanswer}}
\end{yourturn}

% =====================================================================
\section*{\sffamily Mixed set \textbullet\ twelve questions, no labels}

\begin{apcnote}[Why this section exists]
Every question so far arrived with its ladder attached, so you always
knew which method to use. On the exam you will not. These twelve are
deliberately unlabelled and out of order --- the first decision, every
time, is \emph{which kind of problem is this?}
\end{apcnote}

\begin{yourturn}[MIXED SET --- first six]
\begin{enumerate}[leftmargin=*,itemsep=0.5em,label=(\alph*)]
  \item Find the pH of \SI{0.020}{\Molar} \ce{HNO3}.
        \workspace[1.8cm]
  \item Find the pH of \SI{0.020}{\Molar} \ce{NaOH}.
        \workspace[1.8cm]
  \item Give the conjugate base of \ce{H2PO4^-} and the conjugate acid
        of \ce{HCO3^-}. \workspace[1.8cm]
  \item Is a solution of \ce{K2CO3} acidic, basic or neutral? Explain
        in one sentence. \workspace[1.8cm]
  \item A weak acid has $K_{a} = \num{4.0e-6}$. Find $K_{b}$ for its
        conjugate base. \workspace[1.8cm]
  \item Which is the stronger acid, \ce{HBrO} or \ce{HBrO3}? Give the
        reason. \workspace[1.8cm]
\end{enumerate}
\selfcheck{(a) 1.70 \quad (b) 12.30 \quad (c) \ce{HPO4^2-} and \ce{H2CO3}
\quad 4. basic \quad 5. \num{2.5e-9} \quad 6. \ce{HBrO3}}
\keyonly{\begin{apcanswer}\textbf{(a)} \ce{HNO3} is strong, so
$[\ce{H+}] = \SI{0.020}{\Molar}$ and pH $= -\log(\num{2.0e-2}) =
\mathbf{1.70}$. \textbf{(b)} \ce{NaOH} is strong with one \ce{OH-} per
formula unit, so pOH $= 1.70$ and pH $= 14.00 - 1.70 = \mathbf{12.30}$.
\textbf{(c)} Conjugate base of \ce{H2PO4^-} is \textbf{\ce{HPO4^2-}};
conjugate acid of \ce{HCO3^-} is \textbf{\ce{H2CO3}}. \textbf{(d)}
\textbf{Basic} --- \ce{K+} is a spectator from the strong base \ce{KOH},
while \ce{CO3^2-} is the conjugate base of the weak acid \ce{HCO3^-} and
takes a proton from water. \textbf{(e)} $K_{b} =
\num{1.0e-14}/\num{4.0e-6} = \mathbf{\num{2.5e-9}}$. \textbf{(f)}
\textbf{\ce{HBrO3}} --- more oxygens withdraw electron density from the
\ce{O-H} bond and delocalise the conjugate base's negative
charge.\end{apcanswer}}
\end{yourturn}

\begin{yourturn}[MIXED SET --- last six]
\begin{enumerate}[leftmargin=*,itemsep=0.5em,label=(\alph*)]
  \item Find the pH of \SI{0.250}{\Molar} of a weak acid with
        $K_{a} = \num{4.0e-5}$. \workspace[2.2cm]
  \item Find the percent ionisation for question 7.
        \workspace[1.8cm]
  \item Find the pH of \SI{0.0050}{\Molar} \ce{Ba(OH)2}.
        \workspace[2.0cm]
  \item A solution has pOH 4.20. Find its pH and $[\ce{H+}]$.
        \workspace[2.0cm]
  \item An acid has p$K_{a} = 7.2$ and sits at pH 5.0. Which form
        dominates? \workspace[1.8cm]
  \item Explain in one sentence why \SI{0.10}{\Molar} \ce{HCl} and
        \SI{0.10}{\Molar} \ce{CH3COOH} have different pH values.
        \workspace[2.0cm]
\end{enumerate}
\selfcheck{(a) 2.50 \quad (b) 1.3\% \quad (c) 12.00 \quad (d) pH 9.80,
\SI{1.6e-10}{\Molar} \quad 1(a) \ce{HA} \quad 1(b) one is strong, one
weak}
\keyonly{\begin{apcanswer}\textbf{(a)} $x^{2}/0.250 \approx
\num{4.0e-5}$, so $x^{2} = \num{1.0e-5}$ and $x =
\num{3.16e-3}$; pH $= \mathbf{2.50}$. (Check: $1.3\% < 5\%$
\checkmark) \textbf{(b)} $\num{3.16e-3}/0.250 \times 100 =
\mathbf{1.3\%}$. \textbf{(c)} $[\ce{OH-}] = 2 \times 0.0050 =
\SI{0.010}{\Molar}$, pOH $= 2.00$, pH $= \mathbf{12.00}$.
\textbf{(d)} pH $= 14.00 - 4.20 = \mathbf{9.80}$; $[\ce{H+}] =
10^{-9.80} = \mathbf{\SI{1.6e-10}{\Molar}}$. \textbf{(e)}
\textbf{\ce{HA}}, the protonated form --- pH 5.0 is \emph{below} the
p$K_{a}$ of 7.2. \textbf{(f)} \ce{HCl} is a \textbf{strong} acid and
ionises completely, so $[\ce{H+}] = \SI{0.10}{\Molar}$, whereas acetic
acid is \textbf{weak} and only about 1\% ionises, giving far less
\ce{H+} and a much higher pH.\end{apcanswer}}
\end{yourturn}

% =====================================================================
\section*{\sffamily Mastery tracker}

Tick a row only if \textbf{all four} YOUR TURN questions were right on
the first attempt. Every row is assessed except the Lewis model inside
Ladder 1 (see its ``For the exam'' note) --- Br\o nsted--Lowry is the
exam's model.

\renewcommand{\arraystretch}{1.25}
\noindent\begin{tabular}{@{}c p{5.7cm} c p{4.9cm}@{}}
\toprule
\textbf{First try?} & \textbf{Skill} & \textbf{Ladder} &
  \textbf{If not, re-read\ldots}\\
\midrule
$\square$ & the three definitions & 1 & the models nest \\
$\square$ & conjugate pairs & 2 & one proton, one charge \\
$\square$ & strong versus weak & 3 & not the same as concentrated \\
$\square$ & $K_{w}$ and autoionisation & 4 & pH 7 only at 25 C \\
$\square$ & the pH scale & 5 & one unit is tenfold \\
$\square$ & pH of a strong acid & 6 & one line, no ICE \\
$\square$ & pH of a strong base & 7 & count the hydroxides \\
$\square$ & $K_{a}$ and weakness & 8 & bigger $K_{a}$, stronger acid \\
$\square$ & pH of a weak acid & 9 & ICE, then take the log \\
$\square$ & percent ionisation & 10 & dilution raises it \\
$\square$ & weak bases and $K_{b}$ & 11 & $x$ is $[\ce{OH-}]$ \\
$\square$ & $K_{a}K_{b} = K_{w}$ & 12 & conjugate pairs only \\
$\square$ & polyprotic acids & 13 & $K_{a1}$ dominates \\
$\square$ & salts & 14 & the weak parent decides \\
$\square$ & structure and strength & 15 & stabilise the conjugate base \\
$\square$ & pH and p$K_{a}$ & 16 & low pH protonates \\
\midrule
$\square$ & \textbf{mixed set} (10 of 12) & --- & whichever it exposed \\
\bottomrule
\end{tabular}

\begin{apcnote}[Scoring yourself honestly]
16/16 and the mixed set: you have the bulk of an 11--15\% unit, and
Chapter 15 --- buffers and titrations --- will feel like a
continuation rather than a new subject.

12--15: look at \emph{which} ladders missed. Ladders 1--8 are
definitions and one-line arithmetic; they repair fast. Ladders 9--14 are
the calculation core, and a miss there is the one to fix now.

Under 12, or a weak mixed set with strong ladders: that specific
pattern means you can execute each method but cannot yet \emph{choose}
between them --- which is what the exam actually tests. The repair is
not more of the same practice; it is to go back through the ladders
asking only ``what tells me this is a weak-acid problem rather than a
strong-acid one?'' Build the decision, not the arithmetic.

\textbf{One thing to carry into Chapter 15.} Every buffer calculation
starts from a weak acid and its conjugate base --- Ladders 9, 12 and 16
combined. If those three are solid, buffers are mostly bookkeeping.
\end{apcnote}

\end{document}
